%\documentclass[12pt,a4paper,oneside]{scrbook} % la clase book del Koma-script bundle
\documentclass[a4paper,10pt,oneside,headinclude,titlepage]{article} % la clase article estándar
%\linespread{1.25}
\usepackage{setspace}
\usepackage{pdfpages}
\usepackage{hyperref}
\hypersetup{colorlinks,citecolor=red,linkcolor=red,urlcolor=blue,linktocpage}
%\usepackage{tikz}
%\usetikzlibrary{fit,shapes}
\usepackage[spanish]{babel}
%\usepackage{verbatim} %para el entorno comment
%\usepackage{moreverb} %para los ejemplos de lilypond, aporta verbatimtabinput
%\usepackage{alltt} %para los ejemplos de lilypond, aporta verbatiminput
%\usepackage{sverb} %para los ejemplos de lilypond, aporta verbinput
%\usepackage{fancyvrb} %para los ejemplos de lilypond, aporta VerbatimInput
%\pagestyle{empty}
\usepackage[utf8]{inputenc}
\usepackage[T1]{fontenc} %posiblemente sirva para eliminar el problema del enguionado de palabras acentuadas. Lo quitamos provisionalmente para evitar un error
\usepackage{textcomp} % recomendación de Javier Bezos para completar la fuente

\usepackage[margin=2cm]{geometry}
\usepackage{graphicx}
%\usepackage{url}

\usepackage[utopia]{mathdesign}
%\usepackage{mathptmx} %mejor que Times    % alternativa a Charter


%\typearea[0mm]{13}% same as class options above
%\usepackage{newcent}
%\addtokomafont{part}{\mdseries} %encabezamientos sin negrita
%\addtokomafont{partnumber}{\mdseries} %encabezamientos sin negrita
%\addtokomafont{chapter}{\mdseries} %encabezamientos sin negrita
%\setkomafont{disposition}{\normalcolor\bfseries} %no sans serif
%\setkomafont{disposition}{\normalcolor\mdseries} %no negrita

\parskip=0pt\clubpenalty=10000\widowpenalty=10000
\parindent=6mm
\newcommand{\preLilyPondExample}{\vspace{0pt}}
\newcommand{\postLilyPondExample}{\vspace{6pt}}

\newcommand{\lpversion}{2.15.30}
\newcommand{\defsep}{\textbf{$\|$}}
\newcommand{\software}{\emph{software}}
\newcommand{\negspace}{\vspace{-10pt}}  %{\vspace{-20pt}}
\newcommand{\seppar}{
\bigskip
%\vspace{6pt}
}



%%%%%%%%%%%%%%%%%%%%%%%%%%%%%%%%%%%%%%%%%% TÍTULOS
\title{QRTest\\QR Code Test Document}
%\subtitle{}
% necesita Koma-script 2.98 2007/12/24
\author{Francisco Vila}
\date{\today}
%\date{}
%%%%%%%%%%%%%%%%%%%%%%%%%%%%%%%%%%%%%%%%%% FIN TÍTULOS

% \includeonly{apuntes-29-espanol}

\begin{document}

\nonfrenchspacing

\begin{titlepage} %%%%%%%%%%%% PORTADA
  \makeatletter
  \begin{center}
    %\textbf{\Large Universidad de Extremadura}\par
    \vfill
    % \includegraphics[width=60mm]{lily-logo.png}\par
    \vfill
    \textbf{\huge\@title}\par
%    \textbf{\large\@subtitle}\par
    {\@date}
    \vfill
    \textbf{\large\@author}
    \vfill
  \end{center}
  \makeatother
\end{titlepage}

%\singlespace

\begin{singlespace} % los listados van a un espacio. Necesita el paquete setspace
  \tableofcontents
 % \listoffigures
 % \listoftables
\end{singlespace}
 
 % \section*{Introducción}

Esto no es un curso completo de LilyPond ni creo que sustituya a unas
clases directas, tan sólo pretende servir como material de apoyo. La
forma de utilizar estos ejercicios es bastante obvia si se examina uno
cualquiera de ellos:

\begin{enumerate}
\item Observar el modelo y buscar los elementos desconocidos.

\item Leer cuidadosamente el texto para aprender a realizar estos elementos.
Habrá que adaptarlos para recrear el modelo.

\item Tratar de tipografiar el modelo exactamente.

\end{enumerate}

El orden de los ejercicios es importante porque siempre los elementos
nuevos, necesarios para el modelo, están explicados en el mismo
apartado.

Dos son los complementos necesarios para aprender a tipografiar música
con LilyPond.  El primero es la documentación oficial que está en
lilypond.org; el segundo es la comunidad de usuarios, que está a su
disposición en http://lists.gnu.org/mailman/listinfo/lilypond-es
(lista en español) y en
http://lists.gnu.org/mailman/listinfo/lilypond-user (en inglés).

El estado actual de este documento es ``borrador incompleto''.  Dado
que los ejercicios son semanales, su número al término de la colección
será de 30, justo el número de semanas de un curso académico.

%Antes de publicar este borrador en una forma definitiva, posiblemente
%haya que emplear fragmentos alternativos de ciertas obras protegidas.

Mi deseo es que esta pequeña recopilación sea de utilidad a alguien en
algún lugar.  Gracias por leerla.



 \section{Meta}
\vfill
\begin{centering}
Download this document in PDF format:
\\[2mm]
\includegraphics{qrcode}
\\[2mm]
\href{http://paconet.org/wiki/index.php?title=QRTest}{Link to metadata, source code and more!}

\end{centering}

 

\end{document}

